\shorttitle{Двомасштабне стохастичне моделювання з агрегованим міжспільнотним зв'язком}
\shortauthors{Y. Kuryliak and M.T.M. Emmerich}

\title[mode=title]{Двомасштабне стохастичне моделювання епідемій у мережах спільнот з агрегованим міжспільнотним зв'язком}

\author[1]{Yulian Kuryliak}[orcid=0000-0002-6600-2637]
\cormark[1]
\ead{yulian.a.kuryliak@lpnu.ua}
\affiliation[1]{organization={Інститут комп'ютерних наук та інформаційних технологій, Національний університет ``Львівська політехніка''},
            addressline={вул. Степана Бандери, 12},
            city={Львів},
            postcode={79000},
            country={Україна}}

\author[2]{Michael T. M. Emmerich}[orcid=0000-0002-7342-2090]
\ead{michael.t.m.emmerich@jyu.fi}
\affiliation[2]{organization={Факультет інформаційних технологій, Університет Ювяскюля},
            addressline={P.O. Box 35 (Agora)},
            city={Ювяскюля},
            postcode={FI-40014},
            country={Фінляндія}}

\cortext[1]{Автор для листування}

\begin{abstract}
У багатомасштабному стохастичному моделюванні детальна динаміка підсистем часто поєднується через агреговані інтерфейси. Ми досліджуємо цю задачу для поширення епідемій у мережах спільнот, де контакти всередині спільнот подано явно, а зв'язки між ними доступні лише в агрегованій формі. Кожна спільнота моделюється рушієм Gillespie, точним умовно на отриманих імпортованих подіях; макрорівень генерує імпорти за агрегованими залежними від стану інтенсивностями, а в базовій реалізації розміщує їх рівномірно серед сприйнятливих вузлів спільноти-одержувача. Ми повністю специфікуємо метод, зокрема його сіткову синхронізацію, і показуємо, що ідеальна версія з неперервним часом точно моделює сурогатну мережу, у якій кожен пучок міжспільнотних ребер замінено рівномірним зв'язком «усі з усіма» з такою самою сумарною вагою. Це подання виявляє чотири джерела розбіжностей: втрату інформації про стан граничних вузлів, рівномірне розміщення імпортів, відсутність вичерпання ребер інтерфейсу та похибку синхронізації. Для випадково розміщених інтерфейсів у задачі з одиничним експортом ми також доводимо, що ідеальний агрегований годинник експорту стохастично випереджає повномережевий. Контрольовані експерименти на синтетичних мережах SI порівнюють метод із повномережевим моделюванням Gillespie у тестах з одиничним експортом і в 48 конфігураціях ланцюгів із десяти спільнот (50--1000 вузлів у спільноті, внутрішні довжини шляхів 1--6, одне або п'ять ребер інтерфейсу, по 400 зіставлених запусків). Коли розкид часів зараження всередині спільноти був меншим за чверть середнього часу переходу через одне ребро, розподіли часу прибуття до кінцевої спільноти були статистично нерозрізненними у 15 із 18 конфігурацій; за більших значень двомасштабний симулятор прибував раніше на величину, приблизно пропорційну цьому відношенню, яка сягала 40\%. Близькі часи не гарантували допустимих маршрутів: до 52\% агрегованих переходів відбувалися до зараження будь-якого справжнього вузла-моста.
\end{abstract}

\begin{highlights}
\item Точні локальні рушії Gillespie зв'язані через агреговані міжспільнотні інтенсивності
\item Ідеальна схема точно моделює сурогатну мережу з інтерфейсом «усі з усіма»
\item Для випадкових інтерфейсів агрегований годинник експорту стохастично випереджає повномережевий
\item Розбіжність упорядковується співвідношенням внутрішнього розкиду часу та часу переходу через ребро
\item Близькі часи прибуття можуть супроводжуватися фізично неможливими маршрутами передавання
\end{highlights}

\begin{keywords}
Стохастичне моделювання \sep Багатомасштабне моделювання \sep Алгоритм Gillespie \sep Агреговані зв'язки \sep Точність моделі \sep Мережі спільнот
\end{keywords}

\maketitle

%======================================================================
\clearpage
\section{Вступ}
\label{sec:introduction}

Стохастичне моделювання епідемій у структурованих популяціях має враховувати динаміку на різних просторових і структурних масштабах. Усередині спільнот передавання залежить від докладної мережі контактів і локальних стохастичних взаємодій \citep{keeling2005,pastorsatorras2015,kiss2017}, тоді як між спільнотами його часто подають агрегованими потоками мобільності або зваженими мережами зв'язків \citep{balcan2009,barbosa2018human}. Така невідповідність рівнів деталізації створює багатомасштабну задачу: локальну динаміку з високою деталізацією потрібно поєднати через грубіше міжспільнотне подання.

На практиці внутрішньоспільнотні мережі контактів можуть бути спостережені, реконструйовані або синтезовані за неповними даними \citep{eubank2004,danon2011}, тоді як взаємодії між спільнотами можуть бути доступні лише як обсяги маятникової міграції, показники мобільності чи інші агреговані величини \citep{barbosa2018human,chang2021}. Агрегована зв'язність між двома спільнотами не визначає однозначно мікроскопічні контакти, що її утворюють. Тому агреговане подання не зберігає інформації про те, які саме люди формують інтерфейс і як їхні стани впливають на передавання між спільнотами.

Якщо повна мережа контактів відома, епідемічну динаміку можна задати як неперервний марковський ланцюг (CTMC) і точно моделювати подієвими методами типу Gillespie \citep{gillespie1977exact,gibson2000}. Розширення для часових, великих і розріджених мереж підвищують обчислювальну ефективність \citep{vestergaard2015,kuryliak2025simulating}, але все одно потребують явного подання відповідних мікроскопічних каналів взаємодії. На протилежному полюсі метапопуляційні моделі описують взаємодії між спільнотами агрегованими потоками, але зазвичай агрегують і внутрішньоспільнотний стан \citep{colizza2007reaction,colizza2008,balcan2010modeling,ajelli2010}. Між цими крайнощами є проміжна задача: зберегти явну стохастичну динаміку мережі всередині спільнот, але зв'язувати їх через інтерфейс, доступний лише в агрегованій формі.

Двомасштабний симулятор для такого проміжного режиму стисло описано в матеріалах семінару \citep{kuryliak2025towards}. Кожна спільнота моделюється незалежним мікросимулятором, а макрорівень генерує міжспільнотні імпорти за залежними від стану агрегованими інтенсивностями й розміщує кожне зараження у спільноті-одержувачі. Така архітектура зберігає докладну стохастичну динаміку всередині спільнот і агрегує лише інтерфейс зв'язку. У матеріалах семінару обґрунтовано архітектуру та описано її рівні, але не виміряно втрату точності через заміну мікроскопічних міжспільнотних каналів агрегованим зв'язком і правилом розміщення імпортів. Тому тут метод подано самодостатньо й на рівні реалізації.

Центральне питання статті: чи може агрегований міжспільнотний інтерфейс із достатньою точністю відтворити стохастичну поведінку відповідного повномережевого CTMC, для яких спостережуваних величин і в яких мережевих режимах? Відповідь залежить від інформації, яку інтерфейс відкидає: які вузли несуть міжспільнотні канали, куди потрапляє імпортоване зараження та скільки разів можна перетнути тонкий інтерфейс. Ми аналізуємо ці ефекти теоретично й у контрольованих експериментах на синтетичних мережах, де повна мікроскопічна мережа відома досліднику, але двомасштабному симулятору передається лише її агреговане подання.

Основні внески цієї роботи такі.
\begin{enumerate}
\item Повна специфікація двомасштабного методу на рівні реалізації, зокрема застосованого на практиці сітково синхронізованого планувальника на основі інтеграла інтенсивності (Розділ~\ref{sec:method}, Алгоритм~\ref{alg:scheduler}).
\item Точна характеристика ідеального методу як повномережевого моделювання на сурогатній мережі з рівномірним міжспільнотним зв'язком «усі з усіма» (Твердження~\ref{prop:surrogate}). Вона дає таксономію розбіжностей і результат щодо стохастичного впорядкування одиничних експортів (Твердження~\ref{prop:ordering}).
\item Контрольоване порівняння з повномережевим моделюванням Gillespie у тестах з одиничним експортом і на ланцюгах спільнот із використанням перестановочних базових рівнів, що відокремлюють похибку апроксимації від вибіркової мінливості (Розділи~\ref{sec:design} і~\ref{sec:results}).
\item Практичний критерій: відношення розкиду часів зараження всередині спільноти до часу переходу через одне ребро впорядковує спостережувані розбіжності для різних розмірів спільнот, внутрішніх структур і ширин інтерфейсу (Розділ~\ref{sec:res_descriptor}).
\end{enumerate}

Дослідження обмежене динамікою «сприйнятливий--інфікований» (SI) на статичних синтетичних мережах і фіксованим кроком синхронізації. У Розділі~\ref{sec:related} метод зіставлено з наявними підходами до моделювання; у Розділі~\ref{sec:method} його специфіковано; у Розділі~\ref{sec:analysis} проаналізовано, що зберігає агрегований інтерфейс; у Розділах~\ref{sec:design} і~\ref{sec:results} описано експерименти та їхні результати; у Розділі~\ref{sec:discussion} обговорено практичне застосування, обчислювальні витрати й обмеження.

%======================================================================
\clearpage
\section{Пов'язані праці}
\label{sec:related}

\paragraph{Точне та наближене стохастичне моделювання.}
Прямий метод Gillespie точно генерує траєкторії CTMC за поточними схильностями реакцій \citep{gillespie1977exact}; метод наступної реакції організовує ті самі обчислення навколо окремих годинників каналів \citep{gibson2000}. Коли схильності змінюються з часом, точне моделювання можна виразити через подання процесу підрахунку випадковою заміною часу \citep{ethier1986}: подія відбувається, коли її інтегрована інтенсивність досягає незалежного експоненційного порога. \citet{anderson2007} перетворив це подання на модифікований метод наступної реакції для залежних від часу схильностей. Макрогодинник двомасштабного симулятора є прикладом такої конструкції зі стохастичною, залежною від стану інтенсивністю. Орієнтовані на мережі реалізації оновлюють лише схильності, на які вплинула подія \citep{vestergaard2015,kuryliak2025simulating}; використані тут локальні рушії відповідають \citet{kuryliak2025simulating}.

Гібридні та багатомасштабні алгоритми стохастичного моделювання пришвидшують системи з істотно розділеними швидкостями реакцій, наближено відтворюючи швидкі канали \citep{haseltine2002,cao2005slow}. Двомасштабний метод відрізняється об'єктом апроксимації: він не усереднює швидкі канали, а вилучає інформацію про те, які міжспільнотні канали активні. Проте, як показано в Розділі~\ref{sec:results}, його точність також упорядковується розділенням часових масштабів, а саме між внутрішнім поширенням і переходом через інтерфейс.

\paragraph{Метапопуляційні моделі та моделі мобільності.}
Метапопуляційні моделі зв'язують субпопуляції потоками мобільності; їх застосовували для дослідження порогів вторгнення та глобального поширення \citep{colizza2007reaction,colizza2008,balcan2009,balcan2010modeling}; мережі мобільності високої роздільної здатності можна поєднувати з локальними моделями передавання \citep{chang2021}. Агентні моделі описують окремих індивідів у всій системі \citep{eubank2004,kerr2021}, а \citet{ajelli2010} порівняв ці дві парадигми. Двомасштабний метод зберігає окремі вузли та контакти всередині спільнот, а отже, локальну стохастичність і мережеву структуру, проте між спільнотами використовує лише інформацію метапопуляційного типу.

\paragraph{Два рівні перемішування та спільне моделювання.}
Епідемічні моделі з двома рівнями перемішування поєднують локальні контакти в малих групах із рівномірним глобальним перемішуванням \citep{ball1997}. Сурогатна мережа з Твердження~\ref{prop:surrogate} має таку саму структуру, де глобальне перемішування обмежене парами спільнот і зважене агрегованим зв'язком. З погляду програмної організації моделювання двомасштабний планувальник є головним компонентом спільного моделювання, який просуває підсистеми як «чорні скриньки» по комунікаційній сітці й обмінюється невеликою кількістю змінних зв'язку \citep{gomes2018cosim}. На відміну від типових інтерфейсів спільного моделювання, тут інтерфейс не лише передає стан, а й визначає, коли і де у підсистемі-одержувачі відбуваються стохастичні події.

%======================================================================
\clearpage
\section{Метод двомасштабного моделювання}
\label{sec:method}

У цьому розділі визначено повномережевий мікроскопічний еталонний процес і двомасштабний симулятор, використані в роботі. Еталонна модель є CTMC на заданій мережі контактів. Двомасштабний симулятор зберігає внутрішньоспільнотні канали переходів, але замінює явні міжспільнотні канали агрегованим стохастичним інтерфейсом. Ми розрізняємо ідеальний двомасштабний процес із неперервним часом і використану в експериментах реалізацію із синхронізацією на часовій сітці.

\subsection{Повномережева еталонна модель}
\label{sec:reference}

Нехай $G=(V,E)$ --- мережа контактів із множиною вузлів, розбитою на $M$ неперетинних спільнот,
\begin{equation}
\label{eq:partition}
V=\bigsqcup_{i=1}^{M}V_i,\qquad N_i=|V_i|.
\end{equation}
Множину $E$ розглядаємо як набір орієнтованих каналів передавання; неорієнтований контакт утворює два канали. Канали поділено на внутрішньоспільнотні пучки $E_{ii}$ і міжспільнотні пучки $E_{ij}=\{(u,v)\in E:u\in V_i,\,v\in V_j\}$, $i\neq j$. Кожен вузол є сприйнятливим ($S$), інфікованим ($I$) або видаленим ($R$). Мікроскопічний стан позначимо $X(t)\in\{S,I,R\}^{V}$, а його обмеження на спільноту $V_i$ --- $X_i(t)=X(t)|_{V_i}$. Визначимо множини сприйнятливих та інфікованих вузлів $S_i(t)=\{v\in V_i:X_v(t)=S\}$ і $I_i(t)=\{v\in V_i:X_v(t)=I\}$; кількості вузлів у них дорівнюють $s_i(t)=|S_i(t)|$ та $y_i(t)=|I_i(t)|$. Нормовані частки становлять $\sbar_i(t)=s_i(t)/N_i$ і $\ibar_i(t)=y_i(t)/N_i$.

Для конфігурації $x\in\{S,I,R\}^V$ позначимо через $x^{v:I}$ і $x^{u:R}$ конфігурації, у яких вузол $v$ став інфікованим або вузол $u$ став видаленим. Повномережевий процес є CTMC із генератором
\begin{equation}
\label{eq:full_generator}
(\mathcal{L}_G f)(x)=\sum_{(u,v)\in E}\beta_{uv}\,\ind\{x_u=I,\,x_v=S\}\bigl[f(x^{v:I})-f(x)\bigr]
+\sum_{u\in V}\gamma\,\ind\{x_u=I\}\bigl[f(x^{u:R})-f(x)\bigr],
\end{equation}
для обмежених тестових функцій $f$. Тут $\beta_{uv}$ --- інтенсивність передавання каналом $(u,v)$, а $\gamma$ --- інтенсивність видалення. Міжспільнотні інтенсивності записуємо як $\beta_{uv}=\beta_{\mathrm{out}}w_{uv}$, де $w_{uv}$ --- вага контакту. Точне моделювання~\eqref{eq:full_generator} позначаємо \Full{}. В усіх експериментах використано динаміку SI ($\gamma=0$).

\subsection{Агрегований інтерфейс}
\label{sec:interface}

Двомасштабний симулятор, позначений \TS{}, отримує внутрішньоспільнотні графи $G_i=(V_i,E_{ii})$ з їхніми інтенсивностями, але не мікроскопічні міжспільнотні пучки. Кожен пучок подано лише його агрегованою вагою
\begin{equation}
\label{eq:weight}
W_{ij}=\sum_{(u,v)\in E_{ij}}w_{uv},\qquad i\neq j,
\end{equation}
яка для одиничних ваг дорівнює кількості каналів із $V_i$ до $V_j$. У прикладних задачах $W_{ij}$ можна оцінити за даними мобільності або контактів, а не обчислювати з відомої мікроскопічної мережі. Макрорівень спостерігає зведення стану спільноти $(s_i,y_i)$ і відповідні частки $\sbar_i=s_i/N_i$ та $\ibar_i=y_i/N_i$. У Таблиці~\ref{tab:information} наведено інформацію, яку використовує кожна модель. На Рисунку~\ref{fig:schematic}(a,b) зіставлено два подання.

\clearpage
\begin{table}[H]
\centering
\caption{Інформація, яку використовують повномережева еталонна модель і двомасштабний симулятор.}
\label{tab:information}
\small
\begin{tabular}{>{\raggedright\arraybackslash}p{0.38\linewidth}>{\raggedright\arraybackslash}p{0.25\linewidth}>{\raggedright\arraybackslash}p{0.27\linewidth}}
\toprule
Елемент & \Full{} & \TS{}\\
\midrule
Внутрішньоспільнотні контакти та інтенсивності & так & так\\
Стани вузлів усередині спільнот & так & так (локальні рушії)\\
Міжспільнотні канали $E_{ij}$ (вузли-мости та кінцеві вузли) & так & ні\\
Агрегована вага $W_{ij}$ & випливає з $E_{ij}$ & так\\
Інтенсивність передавання з $V_i$ до $V_j$ & зі станів каналів у $E_{ij}$ & з $\ibar_i$, $\sbar_j$ і $W_{ij}$\\
Вузол-одержувач передавання & кінцевий вузол каналу, що спрацював & вибирається за правилом розміщення\\
Максимальна кількість передавань через пучок (SI) & кількість каналів & обмежена лише $s_j$\\
\bottomrule
\end{tabular}
\end{table}

\clearpage
\begin{figure}[H]
\centering
\begin{tikzpicture}[
  x=1cm,y=1cm,
  font=\scriptsize,
  >=Latex,
  person/.style={circle,draw=black!60,fill=blue!18,minimum size=4.2mm,inner sep=0pt},
  infected/.style={person,fill=orange!70},
  bridge/.style={person,very thick,draw=black},
  bridgeinf/.style={infected,very thick,draw=black},
  commbox/.style={draw=black!45,dashed,rounded corners=2pt,inner sep=3pt},
  macrobox/.style={draw=black!70,rounded corners=2pt,fill=black!5,align=center,inner sep=3pt},
  link/.style={draw=black!55},
  weak/.style={draw=black!24,line width=.25pt},
  flow/.style={-{Latex[length=1.7mm]},draw=black!65}
]
% (a) Full microscopic interface
\node[font=\bfseries] at (0,2.25) {(a) Full network};
\node[infected] (a1) at (-1.65,1.20) {};
\node[person]   (a2) at (-1.55,0.35) {};
\node[bridgeinf](a3) at (-0.75,0.82) {};
\node[person]   (a4) at (-0.75,1.55) {};
\node[bridge]   (b1) at (0.65,1.15) {};
\node[person]   (b2) at (1.45,1.55) {};
\node[person]   (b3) at (1.55,0.55) {};
\node[bridge]   (b4) at (0.70,0.30) {};
\draw[link] (a1)--(a2)--(a3)--(a1) (a1)--(a4)--(a3);
\draw[link] (b1)--(b2)--(b3)--(b4)--(b1) (b1)--(b3);
\draw[very thick] (a3)--(b1);
\draw[very thick] (a3)--(b4);
\node[commbox,fit=(a1)(a2)(a3)(a4)] {};
\node[commbox,fit=(b1)(b2)(b3)(b4)] {};
\node at (-1.2,-0.10) {$V_i$};
\node at (1.15,-0.10) {$V_j$};

% (b) Two-scale representation
\node[font=\bfseries] at (5.0,2.25) {(b) Two-scale simulator};
\node[macrobox] (macro) at (5.0,1.55) {macro interface\\$W_{ij}$, $(\ibar_i,\sbar_j)$\\$\lambda_{ij}$ + placement};
\node[macrobox,minimum width=2.15cm,minimum height=.70cm] (li) at (3.65,0.35) {local CTMC\\$G_i$};
\node[macrobox,minimum width=2.15cm,minimum height=.70cm] (lj) at (6.35,0.35) {local CTMC\\$G_j$};
\draw[flow] (li.north) -- node[left] {summaries} (macro.south west);
\draw[flow] (macro.south) -- node[right] {import} (lj.north);
\draw[flow] (lj.east) -- ++(0,1.0) -| node[pos=1.08,right] {summaries} (macro.east);
\node at (5.0,-0.35) {node-level dynamics retained locally};

% (c) Surrogate interface
\node[font=\bfseries] at (10.25,2.25) {(c) Surrogate network};
\node[infected] (c1) at (8.85,1.45) {};
\node[person]   (c2) at (8.75,0.55) {};
\node[person]   (c3) at (9.55,1.05) {};
\node[person]   (c4) at (9.45,0.20) {};
\node[person]   (d1) at (11.00,1.45) {};
\node[person]   (d2) at (10.90,0.55) {};
\node[person]   (d3) at (11.70,1.05) {};
\node[person]   (d4) at (11.60,0.20) {};
\foreach \u in {c1,c2,c3,c4}{\foreach \v in {d1,d2,d3,d4}{\draw[weak] (\u)--(\v);}}
\draw[link] (c1)--(c2)--(c4)--(c3)--(c1);
\draw[link] (d1)--(d2)--(d4)--(d3)--(d1);
\node[commbox,fit=(c1)(c2)(c3)(c4)] {};
\node[commbox,fit=(d1)(d2)(d3)(d4)] {};
\node[align=center] at (10.25,-0.35) {$W_{ij}/(N_iN_j)$ per\\cross-community channel};
\end{tikzpicture}
\caption{Схема міжспільнотного інтерфейсу. (a) У повній мережі інфекція переходить лише явними каналами, що примикають до вузлів-мостів. (b) У двомасштабному симуляторі кожна спільнота зберігає локальний рушій Gillespie/CTMC, а макроінтерфейс використовує зведення стану та агреговану вагу $W_{ij}$ для визначення інтенсивності імпорту й розміщення цільового вузла. (c) За рівномірного розміщення ідеальний двомасштабний процес еквівалентний сурогатній мережі з Твердження~\ref{prop:surrogate}, у якій кожна орієнтована міжспільнотна пара вузлів має слабку інтенсивність $\beta_{\mathrm{out}}W_{ij}/(N_iN_j)$. Помаранчеві вузли інфіковані, сині сприйнятливі.}
\label{fig:schematic}
\end{figure}

\subsection{Локальні рушії}
\label{sec:local}

Кожна спільнота моделюється окремим рушієм, який генерує траєкторії CTMC її власного графа $G_i$: зараження вздовж каналів у $E_{ii}$ та видалення інфікованих вузлів зі схильностями з~\eqref{eq:full_generator}, обмеженими на $V_i$. Рушії використовують інкрементну реалізацію Gillespie з \citet{kuryliak2025simulating}, у якій кожен вузол зберігає сумарну вагу своїх контактів з інфікованими сусідами, а оновлюються лише схильності, на які вплинула подія. Кожен рушій має власний незалежний потік псевдовипадкових чисел.

Рушій підтримує дві операції для макрорівня: просування годинника до заданого часу із застосуванням лише подій, що відбулися раніше, та \emph{імпорт} --- переведення заданого сприйнятливого вузла у стан $I$. Імпорт змінює локальний стан та інтенсивності, тому рушій відкидає заплановану подію й генерує нову з оновленого стану. Такий перезапуск статистично коректний для CTMC, оскільки відповідні часи очікування експоненційні. Еквівалентно, реалізація через випадкову заміну часу може зберегти невикористаний бюджет експоненційної інтенсивності після імпорту. За умови повної послідовності часів імпорту й цільових вузлів локальний рушій точно відтворює CTMC внутрішньоспільнотної мережі з цими імпортами як зовнішніми стрибками $S\to I$. Апроксимація виникає у способі генерування імпортів, а не в локальному моделюванні CTMC.

\subsection{Макрорівень}
\label{sec:macro}

Макрорівень розглядає спільноти як вузли зваженого графа з вагами $W_{ij}$. За поточними узагальненнями стану він призначає кожній упорядкованій парі $i\neq j$ інтенсивність
\begin{equation}
\label{eq:macro_rate}
\lambda_{ij}(t)=\beta_{\mathrm{out}}\,W_{ij}\,\ibar_i(t)\,\sbar_j(t),
\qquad
\Theta(t)=\sum_{i\neq j}\lambda_{ij}(t).
\end{equation}
Макроподія означає один імпорт інфекції з однієї спільноти до іншої. За умови макроподії в час $t$, коли $\Theta(t)>0$, пару $(i,j)$ вибирають з імовірністю $\lambda_{ij}(t)/\Theta(t)$, а цільовий вузол $v\in S_j(t)$ --- за правилом розміщення $\Pi_{ij}(\cdot\mid X(t))$. У всій статті використано рівномірне правило $\Pi_{ij}(v\mid X(t))=1/s_j(t)$. За наявності такої інформації можливі правила, зважені за степенем або з урахуванням межі. Реалізація також має мультиплікативний калібрувальний коефіцієнт для~\eqref{eq:macro_rate}; в усіх експериментах його встановлено рівним одиниці.

Рівняння~\eqref{eq:macro_rate} задає ефективну агреговану інтенсивність, а не, загалом, точну інтенсивність реалізованого мікроскопічного пучка. За фіксованих кількостей вузлів у станах вона дорівнює математичному сподіванню інтенсивності пучка, якщо кінцеві вузли міжспільнотних ребер рівномірно розміщені випадково й ми ще не враховуємо їхні конкретні ідентичності. Інтенсивності збігаються для рівномірного зв'язку «усі з усіма» (Наслідок~\ref{cor:exact}), а в інших випадках агрегована величина є наближенням (Розділ~\ref{sec:analysis}). Макрорівень не автономний: $\Theta$ змінюється після локальної події або імпорту, що змінює кількість вузлів у станах. Оскільки між такими подіями зведення стану кусково сталі, $\Theta$ також кусково стала.

\subsection{Час за інтегралом інтенсивності та синхронізація}
\label{sec:timing}

Нехай $t_0$ --- час попередньої макроподії, а $E\sim\Exp(1)$ --- незалежний поріг. В ідеальній конструкції наступна макроподія відбувається в момент
\begin{equation}
\label{eq:hazard_integral}
t^\ast=\inf\Bigl\{t>t_0:\int_{t_0}^{t}\Theta(s)\,ds\ge E\Bigr\},
\end{equation}
За фіксованих траєкторій мікрорівня $\Theta(t)$ є передбачуваною інтенсивністю, а конструкція~\eqref{eq:hazard_integral} відповідає поданню через випадкову заміну часу, або накопичену інтенсивність, \citep{ethier1986,anderson2007}. Безумовно процес макроподій має стохастичну передбачувану інтенсивність, а не детерміновану інтенсивність звичайного неоднорідного пуассонівського процесу. Для сталої $\Theta$ рівняння~\eqref{eq:hazard_integral} зводиться до часу очікування Gillespie $(t^\ast-t_0)=-\ln U/\Theta$, $U\sim\mathcal{U}(0,1)$. У поєднанні з точними локальними рушіями ця побудова визначає \emph{ідеальний} двомасштабний процес, характеристики переходів якого наведено в Розділі~\ref{sec:surrogate}.

Реалізований планувальник (Алгоритм~\ref{alg:scheduler}) трактує локальні рушії як «чорні скриньки», що відкривають свій стан лише після просування. Він ділить час на інтервали синхронізації довжини $\tau$. Позначимо через $A$ наближений накопичений макрорівневий hazard від попередньої макроподії. На кожному (під)інтервалі $[t,t_g)$, де $\Delta=t_g-t$, планувальник оцінює $\Theta$ пробою в середній точці: копії всіх рушіїв, разом зі станами їхніх генераторів випадкових чисел, просуваються до $t+\Delta/2$, де $\Delta=t_g-t$, після чого за копіями обчислюється $\Theta_{\mathrm{mid}}$. Якщо накопичена інтенсивність $A+\Theta_{\mathrm{mid}}\Delta$ залишається нижчою за $E$, фактичні рушії просуваються до $t_g$. Інакше момент перетину інтерполюється як
\begin{equation}
\label{eq:interpolation}
t^\ast=t+\frac{E-A}{\Theta_{\mathrm{mid}}},
\end{equation}
рушії просуваються до $t^\ast$, а пару та цільовий вузол вибирають за їхніми станами в момент $t^\ast$. Потім поріг поновлюється, й у такий самий спосіб обробляється решта інтервалу.

\begin{algorithm}[tbp]
\caption{Сітково синхронізоване двомасштабне моделювання однієї реалізації.}
\label{alg:scheduler}
\small
\begin{tabbing}
\hspace*{1.6em}\=\hspace*{1.4em}\=\hspace*{1.4em}\=\hspace*{1.4em}\=\kill
\textbf{Вхід:} локальні графи $G_i$, ваги $W_{ij}$, інтенсивності, крок $\tau$, горизонт $T_{\mathrm{end}}$, початковий стан.\\[2pt]
1:\> $t\gets 0$;\ \ $A\gets 0$;\ \ згенерувати $E\sim\Exp(1)$\\
2:\> \textbf{while} $t<T_{\mathrm{end}}$ і інфекція ще може поширюватися \textbf{do}\\
3:\>\> $t_g\gets\min(t+\tau,\,T_{\mathrm{end}})$\\
4:\>\> \textbf{while} $t<t_g$ \textbf{do}\\
5:\>\>\> $\Delta\gets t_g-t$; скопіювати кожен рушій разом зі станом генератора випадкових чисел; просунути копії до $t+\Delta/2$\\
6:\>\>\> $\Theta_{\mathrm{mid}}\gets\sum_{i\neq j}\lambda_{ij}$, обчислена за копіями, рівняння~\eqref{eq:macro_rate}\\
7:\>\>\> \textbf{if} $A+\Theta_{\mathrm{mid}}\Delta<E$ \textbf{then} \hspace{1em}{\color{black!60}// до $t_g$ макроподій немає}\\
8:\>\>\>\> просунути всі рушії до $t_g$;\ \ $A\gets A+\Theta_{\mathrm{mid}}\Delta$;\ \ $t\gets t_g$\\
9:\>\>\> \textbf{else}\\
10:\>\>\>\> $t^\ast\gets t+(E-A)/\Theta_{\mathrm{mid}}$;\ \ просунути всі рушії до $t^\ast$\\
11:\>\>\>\> вибрати $(i,j)$ з імовірністю $\lambda_{ij}(t^\ast)/\Theta(t^\ast)$ за станами рушіїв\\
12:\>\>\>\> рівномірно вибрати $v$ з $S_j(t^\ast)$; імпортувати $v$ до рушія $j$ (його заплановану подію відкинути)\\
13:\>\>\>\> $t\gets t^\ast$;\ \ $A\gets 0$;\ \ згенерувати нове $E\sim\Exp(1)$
\end{tabbing}
\end{algorithm}

Для інтерпретації результатів важливі дві властивості планувальника. По-перше, оскільки проби копіюють стани генераторів випадкових чисел, проба відтворює траєкторію, якою відповідний рушій фактично просуватиметься до середньої точки, якщо не відбудеться імпорт. Тому проба в середній точці є квадратурою реалізованої кусково сталої інтенсивності. Вона точна на підінтервалах без локальних подій у жодній спільноті; похибка квадратури та випереджувальний перегляд не більш як на $\tau/2$ зникають при $\tau\to0$. По-друге, якщо локальні рушії повідомляють часи наступних подій, припущення «чорної скриньки» не потрібне. Оскільки $\Theta$ кусково стала між локальними подіями, інтеграл~\eqref{eq:hazard_integral} тоді можна точно обчислювати подієво керованим планувальником, який об'єднує часи наступних подій усіх рушіїв. Це усуває чисельну похибку синхронізації. В експериментах оцінюємо сітковий варіант із кроком $\tau=1$.

\subsection{Обчислювальна структура та передбачене застосування}
\label{sec:cost}

\Full{} зберігає кожен канал, зокрема всі міжспільнотні канали, а його вартість на одну подію зростає зі степенем вузла \citep{kuryliak2025simulating}. \TS{} зберігає внутрішньоспільнотні графи та матрицю ваг розміру $M\times M$. Саме цю економію підкреслено в матеріалах семінару \citep{kuryliak2025towards}: коли міжспільнотні контакти численні, наприклад коли потік мобільності відповідає багатьом індивідуальним контактам, інтерфейс зводиться до $O(M^2)$ чисел. Локальні рушії також можна змінювати незалежно, що спрощує подання втручань, специфічних для окремих спільнот. Водночас сітково синхронізований планувальник додає витрати, пропорційні $(T_{\mathrm{end}}/\tau)\sum_iN_i$, на копіювання станів рушіїв, незалежно від кількості подій. Отже, метод призначений для умов, у яких міжспільнотні контакти невідомі або надто численні для явного подання, а не для малих мереж із доступною повною структурою контактів. Виміряний час виконання наведено в Розділі~\ref{sec:disc_cost}.

%======================================================================
